- citation philosophy
- description vs analysis
- latex -> html
  - Citations can pop-up info box
    - Additional information in citation without inflating citation count
  - Code samples can be in margin
  -references to simple call can pop up the relevant code snippet
- doubts about black box
- doubts about simple call
  - Some language influencing Python used named values as the "lower level" way (modula?)
- thoughts on synthesis
- Items for next draft
  - conversion to HTML
  - first pass at synthesis
- Splitting into *-lib to avoid call-label

There are 2 primary additions contained in this second draft: an expanded introduction
which now includes visualizations to explain how I am reasoning about arguments and the
newly added section on guile.
In the process of adding these things, I have gathered some doubts about how I am
approaching this paper.

The biggest concern I have is the treatment of intermediate representations as black
boxes.
I restrict myself to examining the source code as input and whatever representation the
language considers to be the "final" output - for Python and Guile this is bytecode, but
for C this will be machine code.
This seems useful to contain the scope of the paper, as otherwise the analysis section
could get out of hand (if it is not already).
However, I am worried that ignoring the intermediate representations prevents me from
attaining certain insights.
The purpose of examining the output of compilers is to understand how the computer is
processing the information given to it by the human and to inform the implementation of
the framework.
The end result gives a concrete and literal understanding of what actions the computer
decides to take based on the human's input, but it doesn't tell me how the computer
arrived at those conclusions or why it made the decisions it did.
What happened during the Guile section is that I looked at intermediate representations
in order to inform my description of the language, but I did not directly discuss these
representations in the paper.
This might be the best compromise, and it is how I will proceed for the time being.

I am uncertain what subject of study this paper actually falls into.
Looking at the paper by word count, much of it is focused on the way that programming
languages operate and accept input - so, computer science right?
Except that my primary concern is not about the compuer per se, but the way that humans
use source code to communicate, both with the computer and with each other.
So is this information science because I'm concerned with the transmission of information?
Is it sociological because I'm concerned with human-to-human interaction, albeit indirect
and asynchronous interaction?
Or is it just computer science that (should) draw on other fields in order to enhance the
validity of the work?
I'm not sure.
I'm also uncertain how useful it is to precisely define what the "subject of study" is, or
if it is even realistic to expect that any given work will fit neatly into exactly one
subject of study.
